\documentclass[11pt,a4paper]{article}
\usepackage[margin=2.1cm]{geometry}
\usepackage[T1]{fontenc}
\usepackage[utf8]{inputenc}
\usepackage[brazilian]{babel}
\usepackage{lmodern}
\DeclareUnicodeCharacter{2014}{\textemdash}
\DeclareUnicodeCharacter{2013}{\textendash}
\DeclareUnicodeCharacter{2022}{\textbullet}
\usepackage{amsmath,amssymb,booktabs,array,xcolor,hyperref,fvextra,upquote}
\hypersetup{colorlinks=true,linkcolor=blue,urlcolor=blue}
\fvset{breaklines=true,breakanywhere=true}
\sloppy
\definecolor{darkblue}{HTML}{193A5E}
\newcommand{\code}[1]{\texttt{#1}}
\newcommand{\commandblock}[1]{\begin{quote}\small\ttfamily #1\end{quote}}
\setlength{\parskip}{0.6em}
\setlength{\parindent}{0pt}
\begin{document}

\begin{center}
{\LARGE\bfseries\color{darkblue} Relatório de Criptoanálise --- Grupo 7}\\[0.6em]
{\large Recuperação completa da mensagem confidencial}\\[0.5em]
Luísa Carvalho Böhm \textbullet{} Murilo Costa Salem \textbullet{} Henrique dos Reis\\
25 de outubro de 2026
\end{center}

\begin{abstract}
Este relatório descreve, na ordem em que foram executadas, a descoberta da senha de acesso à chave pública, a decifragem AES dessa chave, a fatoração do módulo RSA, a reconstrução da chave privada, a recuperação da senha da mensagem e a decifragem do texto final. Inclui parâmetros, comandos reproduzíveis, tentativas malsucedidas e o código integral do script desenvolvido para a reconstrução RSA.
\end{abstract}

\tableofcontents
\newpage

\section{Arquivos de entrada e ambiente}

Os quatro arquivos fornecidos foram \code{key\_for\_rsa\_public.hash}, \code{key\_public.en}, \code{key\_for\_message.en} e \code{message.en}, todos no diretório \code{7/}. O enunciado informou que a senha do primeiro arquivo começava por \code{pe}. Os comandos abaixo partem da raiz do projeto, em Linux, com Hashcat 7.1.2, OpenSSL 3.5.5, Python 3 e a biblioteca Python \code{cryptography} 46.0.5. A busca com Hashcat utilizou uma NVIDIA GeForce RTX 5090, identificada como dispositivo 1 via OpenCL.

Se o Hashcat ainda não estiver no projeto, obtenha e extraia a mesma versão antes da etapa de força bruta:
\begin{Verbatim}[fontsize=\scriptsize]
mkdir -p .tools
curl -fL -o .tools/hashcat-7.1.2.7z \
  https://github.com/hashcat/hashcat/releases/download/v7.1.2/hashcat-7.1.2.7z
7z x -y .tools/hashcat-7.1.2.7z -o.tools
\end{Verbatim}
Para reconstruir a chave RSA em Python, instale \code{cryptography} se necessário. As demais ferramentas usadas nos comandos são OpenSSL, Docker, \code{xxd} e \code{7z}.

\begin{center}
\begin{tabular}{p{0.30\linewidth}p{0.61\linewidth}}
\toprule
Arquivo & Formato identificado \\
\midrule
\code{key\_for\_rsa\_public.hash} & SHA-256 em hexadecimal, 64 caracteres \\
\code{key\_public.en} & Base64 de ciphertext AES, sem cabeçalho \code{Salted\_\_} \\
\code{key\_for\_message.en} & Ciphertext RSA em hexadecimal, 48 bytes após conversão \\
\code{message.en} & Base64 do formato \code{openssl enc} com cabeçalho \code{Salted\_\_} \\
\bottomrule
\end{tabular}
\end{center}

\section{Descoberta da senha da chave pública}

O hash informado é:
\begin{Verbatim}[fontsize=\scriptsize]
bfc15bd9ea199751e66ddb5e92344baed2f2fe67136e12eea18281ca52ff61fb
\end{Verbatim}

O modo 1400 do Hashcat corresponde a SHA-256 bruto; \code{-a 3} seleciona busca por máscara e \code{?l} representa letras minúsculas. A tentativa inicial exigiu exatamente seis letras minúsculas após \code{pe}:
\begin{Verbatim}[fontsize=\scriptsize]
./.tools/hashcat-7.1.2/hashcat.bin -m 1400 -a 3 \
  7/key_for_rsa_public.hash 'pe?l?l?l?l?l?l' -O -w 3 -d 1 \
  --potfile-disable --restore-disable
\end{Verbatim}
Ela terminou com \code{Status: Exhausted}: os $26^6=308.915.776$ candidatos foram testados sem correspondência. A máscara foi então ampliada para dez letras minúsculas após \code{pe}, ou 12 caracteres no total:
\begin{Verbatim}[fontsize=\scriptsize]
./.tools/hashcat-7.1.2/hashcat.bin -m 1400 -a 3 \
  7/key_for_rsa_public.hash 'pe?l?l?l?l?l?l?l?l?l?l' \
  -O -w 3 -d 1 --potfile-disable --restore-disable
\end{Verbatim}
O resultado foi \textbf{\code{petershorxxx}}. O espaço máximo dessa segunda máscara é $26^{10}=141.167.095.653.376$ candidatos. O Hashcat informou \code{Status: Cracked} após cerca de 12 segundos de ataque, com 193.476.689.920 candidatos contabilizados (0,14\% do espaço). Para conferir diretamente a senha descoberta:
\begin{Verbatim}[fontsize=\scriptsize]
printf %s petershorxxx | openssl dgst -sha256
\end{Verbatim}
O digest deve coincidir com o conteúdo de \code{key\_for\_rsa\_public.hash}.

\section{Decifragem de \code{key\_public.en}}

O conteúdo Base64 desse arquivo não começa com o cabeçalho \code{Salted\_\_}, de modo que a decifragem OpenSSL precisa da opção \code{-nosalt}. A combinação validada foi \textbf{AES-256-OFB}, codificação Base64, derivação legada de senha do OpenSSL com digest \textbf{SHA-256}, sem salt:
\begin{Verbatim}[fontsize=\scriptsize]
openssl enc -d -aes-256-ofb -a -nosalt -md sha256 \
  -pass pass:petershorxxx -in 7/key_public.en -out 7/key_public.pem
openssl pkey -pubin -in 7/key_public.pem -noout -text
\end{Verbatim}
O resultado é uma chave pública RSA de 382 bits, expoente $e=65.537$, em formato PEM \code{BEGIN PUBLIC KEY}. O módulo decimal é:
\begin{Verbatim}[fontsize=\tiny]
6524226763420669545131292616008021977421989716443706623412713108662980402279261528264966716400244521958486921252927
\end{Verbatim}
O SHA-256 da chave pública no formato DER é:
\begin{Verbatim}[fontsize=\scriptsize]
dbf1517e0312436f41a2b683b296a065937e868e72a65b7a7dfc445bff32f961
\end{Verbatim}
O aviso do OpenSSL sobre derivação obsoleta é esperado: esses arquivos foram produzidos com essa derivação; usar PBKDF2 mudaria a chave gerada.

\section{Fatoração do módulo e reconstrução da chave privada}

A ferramenta CADO-NFS 3.0.0, com o crivo geral de corpos numéricos (GNFS), fatorou o módulo. Para reproduzir com 20 threads e estado persistido em um diretório local:
\begin{Verbatim}[fontsize=\scriptsize]
mkdir -p .factor-rsa382
RSA_N='652422676342066954513129261600802197742198971644370662341271'
RSA_N="${RSA_N}3108662980402279261528264966716400244521958486921252927"
docker run --name factor-rsa382 \
  -v "$PWD/.factor-rsa382:/work" \
  -e RSA_N="$RSA_N" \
  registry.gitlab.inria.fr/cado-nfs/cado-nfs/factoring-full:latest \
  /bin/bash -lc 'cado-nfs.py --workdir /work "$RSA_N" -t 20'
\end{Verbatim}
O tempo de parede registrado foi 935,176 segundos. Os fatores primos obtidos foram:
\begin{Verbatim}[fontsize=\scriptsize]
p = 2295796704686193901538758441276290741108975069288091220567
q = 2841813802634779914926926270734337524686850991346453461081
\end{Verbatim}
O produto $pq$ foi comparado com o módulo da chave pública, e ambos os fatores foram verificados como primos pelo OpenSSL. A partir daí foi calculado $\lambda(n)=\operatorname{mmc}(p-1,q-1)$ e o expoente privado $d\equiv e^{-1}\pmod{\lambda(n)}$. O script integral está no Apêndice B. Para executá-lo em uma cópia limpa dos arquivos (sem as saídas pré-existentes):
\begin{Verbatim}[fontsize=\scriptsize]
python3 7/recover_rsa.py
openssl pkey -in 7/key_private.pem -check -noout
\end{Verbatim}
O script grava \code{key\_private.pem} em PEM PKCS\#8 sem senha e com permissão 0600. O OpenSSL retornou \code{Key is valid}; a chave pública extraída da privada tem o mesmo SHA-256 DER indicado acima. A chave privada é um resultado sensível: mantenha suas permissões restritas.

\section{Recuperação da senha da mensagem}

O arquivo \code{key\_for\_message.en} contém 96 caracteres hexadecimais, que representam um bloco RSA de 48 bytes. O script converte o hexadecimal em bytes e aplica a operação privada com padding \textbf{PKCS\#1 v1.5}. A estrutura interna observada começa por \code{00 02}, contém bytes de preenchimento não nulos e um separador \code{00}. Após esse separador, o texto recuperado é:
\begin{center}
\Large\bfseries\code{kristennygaa}
\end{center}
O valor foi salvo em \code{7/key\_for\_message.txt}, sem quebra de linha. Este comando OpenSSL reproduz a decifragem a partir da chave privada já criada:
\begin{Verbatim}[fontsize=\scriptsize]
tr -d '\r\n' < 7/key_for_message.en | xxd -r -p > /tmp/key_for_message.bin
openssl pkeyutl -decrypt -inkey 7/key_private.pem \
  -in /tmp/key_for_message.bin -out /tmp/key_for_message_recovered.txt
\end{Verbatim}
Por padrão, a operação RSA acima usa PKCS\#1 v1.5. Compare o conteúdo do arquivo recuperado em \code{/tmp/} com \code{kristennygaa}.

\section{Decifragem da mensagem final}

Após conversão Base64, \code{message.en} começa com \code{Salted\_\_}, cabeçalho do formato \code{openssl enc} que inclui um salt aleatório. Testes com AES e diferentes derivações produziram um único texto UTF-8 inteiramente legível: \textbf{AES-256-OFB}, Base64, salt embutido, derivação legada OpenSSL com digest \textbf{SHA-256}. O comando é:
\begin{Verbatim}[fontsize=\scriptsize]
openssl enc -d -aes-256-ofb -a -md sha256 \
  -pass file:7/key_for_message.txt \
  -in 7/message.en -out 7/message.dec.txt
\end{Verbatim}
O texto final contém exatamente os nomes Luísa Carvalho Böhm, Murilo Costa Salem e Henrique dos Reis, conforme exigido. O SHA-256 do arquivo UTF-8 recuperado é:
\begin{Verbatim}[fontsize=\scriptsize]
674d3ec5a364f98594e4edf2684680fc348e76eefa78ab8185b9b772d6148b7e
\end{Verbatim}
O Apêndice A contém a mensagem integral, com todas as linhas preservadas.

\section{Caminhos testados que não funcionaram}

\begin{enumerate}
\item \textbf{Máscara com seis letras após \code{pe}.} O Hashcat testou todo o espaço de 308.915.776 possibilidades e informou \code{Exhausted}. Falhou porque a senha real possui dez letras após \code{pe}; nenhuma senha de oito caracteres naquela máscara poderia produzir o hash alvo.
\item \textbf{Aplicação da chave pública ao bloco RSA.} A tentativa \code{openssl pkeyutl -verifyrecover -pubin} com padding PKCS\#1 v1.5 retornou \code{invalid padding}. Usando a operação pública sem padding, o resultado foi um bloco binário sem a senha. Tecnicamente, o ciphertext deste arquivo foi recuperado pela operação com o expoente \emph{privado} e padding PKCS\#1 v1.5; a operação pública aplica o expoente errado para essa inversão.
\item \textbf{Decifragem de \code{key\_public.en} como arquivo OpenSSL com salt.} A tentativa com \code{openssl enc -d -aes-256-cbc -a -pbkdf2} retornou \code{bad magic number}: os bytes decodificados não começam por \code{Salted\_\_}. Esse erro evidencia a ausência do cabeçalho esperado; o modo AES, a derivação e o uso de salt também não coincidem com os parâmetros corretos.
\end{enumerate}

\section{Resumo dos parâmetros encontrados}
\begin{center}
\begin{tabular}{p{0.30\linewidth}p{0.61\linewidth}}
\toprule
Camada & Resultado \\
\midrule
Senha de \code{key\_public.en} & \code{petershorxxx}; SHA-256 bruto no arquivo de hash \\
\code{key\_public.en} & AES-256-OFB; Base64; sem salt; KDF legado do OpenSSL com SHA-256 \\
\code{key\_for\_message.en} & RSA de 382 bits; ciphertext hexadecimal; padding PKCS\#1 v1.5 \\
Senha recuperada para a mensagem & \code{kristennygaa}, 12 bytes ASCII \\
\code{message.en} & AES-256-OFB; Base64; salt no cabeçalho \code{Salted\_\_}; KDF legado do OpenSSL com SHA-256 \\
Texto final & UTF-8; reproduzido integralmente no Apêndice A \\
\bottomrule
\end{tabular}
\end{center}

\appendix
\section{Mensagem decifrada completa}
\VerbatimInput[fontsize=\scriptsize]{message.dec.txt}

\section{Código-fonte integral do script RSA desenvolvido}
O arquivo abaixo é o mesmo \code{recover\_rsa.py} utilizado nas etapas 4 e 5. O projeto também contém ferramentas didáticas preexistentes, como \code{hashcat\_python.py} e \code{dictionary\_rules.c}; elas não foram executadas para obter os resultados apresentados neste relatório.
\VerbatimInput[fontsize=\scriptsize]{recover_rsa.py}

\end{document}
